\section{Discussion}\label{sec:discuss}

Pseudo-Hermitian theories provide an alternative framework to define consistent quantum mechanical systems. It might be possible and very interesting to find pseudo-Hermitian formulations of quantum gravity by taking advantage of the holographic duality. We have taken a step in this direction by constructing solutions that can be interpreted as dual to pseudo-Hermitian theories in analytic extensions of well-known consistent supergravity truncations. The fact that we have not found linear instabilities in wormhole solutions despite the presence of ghost scalars is encouraging, although we cannot rule out non-perturbative instabilities, or instabilities in fields not belonging to the consistent truncations. 

In \cite{Maldacena:2026jqd} (see also \cite{DiUbaldo:2026rly,Etheredge:2026rio}) it has been argued that the distance along the imaginary direction that the scalar travels when traversing a wormhole places a bound on the analytic extension of a coupling in the dual field theory. The lower and upper bounds that were found for asymptotically $\text{AdS}_{d+1}$ boundaries were
\begin{equation}
	\IDB_{\text{lower}}(d)
	=
	\pi \sqrt{\frac{2(d-1)}{d}}
	\ , \qquad 
	\IDB_{\text{upper}}(d)
	=
	\pi \sqrt{\frac{2d}{d-1}}\ .
\end{equation}
These bounds were found for massless scalar fields, that are dual to marginal operators. In the examples we have constructed, the scalar operators have a non-trivial potential and the dual operators are relevant. We find that in these cases the distance between the two sides of the wormhole is larger than the upper bounds of the marginal case
\begin{align}
\text{Example 1:}\quad& \Delta \eta=% generated by calc/mathematica -- do not edit
\sqrt{3} \pi
> \IDB_{\text{upper}}(4)=% generated by calc/mathematica -- do not edit
2 \sqrt{\frac{2}{3}} \pi

\ ,\\
\text{Example 2:}\quad& \Delta \eta=% generated by calc/mathematica -- do not edit
2 \pi
> \IDB_{\text{upper}}(3)=% generated by calc/mathematica -- do not edit
\sqrt{3} \pi

\ ,\\
\text{Example 5:}\quad& |\Delta \vec \eta|= \sqrt{(\sqrt{3}\pi)^2+\pi^2} = % generated by calc/mathematica -- do not edit
2 \pi
 > \IDB_{\text{upper}}(4)=% generated by calc/mathematica -- do not edit
2 \sqrt{\frac{2}{3}} \pi
\ .
\end{align}

Concerning fluctuations in wormhole solutions, we find that the spectrum of ghost fields consists of two identical copies, each of them localized on one side of the wormhole. This implies that the two-point correlators of the dual scalar operators do not have mixed contributions between the two QFTs. In fact, there are no double-trace deformations coupling the two copies of the scalar operator, which in other setups are needed to make wormholes traversable \eg \cite{Gao:2016bin,Maldacena:2017axo,Maldacena:2018lmt,Caceres:2018ehr,Bak:2018txn,Ahn:2022suv,Khairunnisa:2025ljk,Djogama:2026pmd}. Fluctuations of the metric, on the other hand, split according to their parity under reflection along the holographic direction, which exchanges the two boudaries. In the dual theory, parity-even and parity-odd modes would be naturally interpreted as diagonal and anti-diagonal combinations of the two copies of the energy-momentum tensor.

The wormhole solutions we have found are Lorentzian, traversable, with boundaries that have flat metrics and no non-local interactions between them. In our interpretation of the solutions, the two boundaries do not correspond to two disconnected spacetimes, but to two decoupled copies of a theory in the same spacetime. The wormhole corresponds to an entangled state of the product Hilbert space. The state spontaneously breaks the two copies of Poincar\'e symmetry to their diagonal subgroup. If the origin of the breaking is tachyon condensation, the diffeomorphism Wilson line operator connecting the two boundaries would correspond to the components of the tachyon along flat directions, acting as the Goldstone bosons of the broken Poincar\'e symmetry. Extending the argument to other symmetries could serve as an alternative resolution--at the level of effective theory--of the factorization problem addressed in \cite{Harlow:2015lma,Guica:2015zpf}. Interestingly, in addition to the natural extension to Euclidean wormholes \cite{Betzios:2023obs}, spontaneous symmetry breaking due to entanglement between two boundaries of a holographic spacetime may be realized more generally. A similar mechanism seems to be at play in holographic entanglement island constructions, where the would-be Goldstone is eaten by a dynamical gauge field at one of the boundaries \cite{Geng:2025byh,Geng:2025gqu,Geng:2025gns,Geng:2025rov}.

Although we have focused on wormhole solutions, some of the other solutions that we have found have interesting characteristics, too. In particular, the flows to asymptotically flat spacetime have an uplift to ten-dimensional singular solutions with a real metric and imaginary fluxes. Example 6 of analytic extension of a supergravity solution, dual to a 3+1 dimensional QFT, is uplifted to a type IIB solution with purely imaginary two-form NS-NS and R-R  potentials. It is interesting to compare with the classification of \cite{Bergshoeff:2007cg}, which is summarized in their table 1. In addition to type IIB, they identify two other theories with real action, IIB${}^*$ and $\text{IIB}^\prime$ (related to each other through S-duality). In both cases, one of the two-form potentials is real and the other is imaginary, so the solution we have found does not belong to either. It might be that there is a more general classification if one relaxes some of the conditions, such as maximal supersymmetry. Example 7 is similar. It is the analytic extension of a supergravity solution dual to a 2+1 dimensional QFT. The uplift is to an eleven dimensional supergraviy solution with four form potential that has both real and purely imaginary components. This falls outside the classification of M-theory solutions given in Appendix D of \cite{Bergshoeff:2007cg}, where all the components would be either real or imaginary.

An interesting extension of this work would be to construct solutions interpolating between anti-de Sitter and de Sitter spacetimes. This could be achieved if there are flows that interpolate between critical points of the superpotential with negative and positive values. Recent works show that this is not possible quite generically if the null energy condition is satisfied \cite{Kiritsis:2019wyk,Kiritsis:2025ytb}. Since the analytic extension of the scalars introduces ghosts that violate the null energy conditions, it is possible that one can go around this obstruction.

